\section{Perturbation Semantics}
\label{sec:perturbation-semantics}

Stability is not defined until both the perturbation operation and the
evaluation point are specified.  This section fixes the notation and
distinguishes four perturbation notions that are often conflated.

\subsection{Ordered finite samples and deterministic $k$-NN}

\begin{definition}[Sample]
A \emph{sample} $S$ of size $n$ is an ordered tuple
\[
S = \big((x_1, y_1), \ldots, (x_n, y_n)\big)
\]
where each $x_i \in X$ lies in a finite metric space $(X, d)$ and each label
$y_i \in \{0, 1\}$.
\end{definition}

The order of the tuple is semantically relevant only for deterministic
tie-breaking.  Two identical labeled points appearing at different indices are
treated as distinct sample occurrences.

\begin{note}[Duplicate points and conflicting labels]
Duplicate points are allowed.  Conflicting labels at the same point are
allowed.  Any theorem or experiment that excludes duplicates or conflicts must
state this explicitly.
\end{note}

These conventions are not cosmetic: duplicates, conflicting labels, and
deterministic ordering are precisely where local perturbation sensitivity can
appear.  Several examples collapse or become ambiguous if these conventions are
left implicit.

\begin{definition}[Finite Metric Space]
A finite metric space is a finite set $X$ together with a distance function
\[
d : X \times X \to \mathbb{R}_{\ge 0}
\]
satisfying for all $x, x', z \in X$: $d(x, x) = 0$; $d(x, x') > 0$ for
$x \ne x'$; $d(x, x') = d(x', x)$; and the triangle inequality.
\end{definition}

\begin{definition}[Graph Shortest-Path Metric]
For a finite, simple, connected, undirected graph $G = (V, E)$, the graph
metric is $d_G(u, v) = \text{length of a shortest path from } u \text{ to } v$.
\end{definition}

\begin{definition}[Deterministic Tie-Breaking]
Neighbor ordering is determined lexicographically by: (1) smaller distance to
the query point; (2) smaller sample index in the ordered tuple.  If class votes
tie after the first $k$ ordered neighbors are selected, the predicted label is
resolved by fixed label order $0 \prec 1$.
\end{definition}

\begin{definition}[$k$-NN Prediction]
\label{def:knn-prediction}
Given $k \in \mathbb{N}$ with $1 \le k \le n$, the deterministic $k$-NN
predictor $h_S^{(k)}(x)$ is obtained by: ordering the sample occurrences by the
deterministic neighbor rule above; selecting the first $k$ occurrences; taking
the majority vote of their labels; resolving label-vote ties by $0 \prec 1$.
\end{definition}

\begin{definition}[Binary 0-1 Loss]
For a hypothesis $h$ and a labeled point $(x, y)$, the binary classification
loss is $\ell(h, (x, y)) = \mathbf{1}\{h(x) \ne y\}$.
\end{definition}

\subsection{Delete-one, insert-one, replace-one}

Let $S = \big((x_1, y_1), \ldots, (x_n, y_n)\big)$.

\begin{definition}[Delete-One]
For index $i \in \{1,\ldots,n\}$, define $S^{-i}$ as the ordered tuple with the
$i$-th occurrence removed.
\end{definition}

\begin{definition}[Insert-One]
For any labeled point $z = (x, y) \in X \times \{0,1\}$ and any insertion
position $j \in \{1,\ldots,n+1\}$, define $S \oplus_j z$ as the ordered tuple
obtained by inserting $z$ at position $j$.
\end{definition}

\begin{definition}[Replace-One]
For index $i$ and labeled point $z = (x', y') \in X \times \{0,1\}$, define
$S^{i \leftarrow z}$ as the ordered tuple obtained by replacing the $i$-th
occurrence of $S$ with $z$.
\end{definition}

\subsection{Why LOO is not delete-one}

\begin{definition}[LOO / CVloo Stability]
For each index $i$, define
\[
\Delta_{\mathrm{loo}}(S, i)
=
\left| \ell(h_{S^{-i}}^{(k')}, (x_i, y_i)) - \ell(h_S^{(k)}, (x_i, y_i)) \right|,
\]
where $k' = \min(k, n-1)$ when deletion reduces sample size below $k$.
\end{definition}

LOO is not simply delete-one.  It is delete-one with the evaluation point
constrained to be the deleted occurrence.  This coupling between perturbation
and evaluation point is the defining property that distinguishes LOO from
replacement-based stability notions.

The following table makes the distinctions explicit.

\begin{table}[t]
\centering
\small
\begin{tabularx}{\linewidth}{@{}lXXl@{}}
\toprule
Notion & Perturbation & Evaluation point & Tests what \\
\midrule
Delete-one & Remove one occurrence & Arbitrary query & Deletion sensitivity \\
Insert-one & Add one occurrence & Arbitrary query & Insertion sensitivity \\
Replace-one & Swap one occurrence & Arbitrary query & Replacement robustness \\
LOO & Delete one occurrence & Deleted occurrence & CV sensitivity \\
\bottomrule
\end{tabularx}
\caption{Four perturbation notions distinguished by what is perturbed and where
the perturbed predictor is evaluated.  The comparison shows that LOO is a
special case of delete-one with a restricted evaluation point.}
\label{tab:perturbation-semantics}
\end{table}

The pointwise stability indicators used throughout the paper are:

\begin{align*}
\Delta_{\mathrm{del}}(S, i, x, y) &=
\left| \ell(h_S^{(k)}, (x, y)) - \ell(h_{S^{-i}}^{(k')}, (x, y)) \right|,
\\
\Delta_{\mathrm{rep}}(S, i, z, x, y) &=
\left| \ell(h_S^{(k)}, (x, y)) - \ell(h_{S^{i \leftarrow z}}^{(k)}, (x, y)) \right|,
\\
\Delta_{\mathrm{ins}}(S, j, z, x, y) &=
\left| \ell(h_S^{(k)}, (x, y)) - \ell(h_{S \oplus_j z}^{(k)}, (x, y)) \right|,
\\
\Delta_{\mathrm{loo}}(S, i) &=
\left| \ell(h_{S^{-i}}^{(k')}, (x_i, y_i)) - \ell(h_S^{(k)}, (x_i, y_i)) \right|.
\end{align*}

For LOO, the evaluation pair $(x_i, y_i)$ is fixed to the deleted occurrence,
so it is not a free parameter.  This is the critical difference: replace-one
and delete-one can be evaluated at any query, while LOO can only testify about
the deleted point itself.
